\documentclass[11pt,a4paper]{article}
\usepackage[margin=25mm]{geometry}
\usepackage{amsmath,amssymb,amsthm}
\setlength{\emergencystretch}{2em}
\newtheorem{lemma}{Lemma}
\newtheorem{theorem}[lemma]{Theorem}
\newcommand{\gen}[1]{\langle#1\rangle}
\title{The location of the exceptional exit in a later shared crossing}
\author{Human Compression project: L044 continuation}
\date{}
\begin{document}\maketitle
Retain the complete standing assumptions and smaller-degree alternating
induction of \emph{Completing the transitive-tail alternative} and
\emph{Routing a maximal shared crossing}. In particular $G\le A_n$
is transitive with a connected string C-marking, every maximal
omission is intransitive, and eligibility is always ambient.
All full predecessor proofs and their required finite inputs remain
acquired dependencies.

\begin{lemma}[The two exits have different index locations]\label{location}
In the exact transitive-tail interface, let $i$ be the unique crossing,
$h>i$ its last shared label and $G_{>h}$ transitive on its support.
Under rank failure, either the actual last three generators form a
six-point Table 2 head, with $h=r-4$, or $h=2$ and the actual first
three generators form such a head. In the latter case $i=0$ or $1$.

In the fully reversed interface, let $j$ be the crossing and $g<j$
its first shared label, with the prefix $G_{<g}$ transitive on its
support. The alternatives are an actual first head with $g=3$, or
an actual last head with $g=r-3$.
\end{lemma}
\begin{proof}
We inspect the full proof, not just its weaker either-end statement.
Its small transitive tail is either the degree-five dihedral group,
which the exact critical-tail lemma excludes, or an actual six-point
$A_5$ marking. The latter has rank three, so $r-1-h=3$.
The strict proper primitive and anchored-imprimitive bounds exclude
all other proper tails.

The remaining natural alternating tail has the precise equality
parameters in equation (2) of the acquired proof. In its endpoint
case the exact fixed-path lemma forces $i=0,h=2$ and the first
head. Its transitive-left interior case is impossible. Its
intransitive-left interior case forces $h=i+1=2$; the complete
six-choice dihedral attachment then gives the first head.
There are no other exits from that proof. Replacing every index
$k$ by $r-1-k$ gives the asserted prefix alternatives: the reversed
last shared index is $r-1-g$, so its values $r-4$ and $2$ become
$g=3$ and $g=r-3$, respectively. This is full string reversal,
not a local reorder of the marked tuple.
\end{proof}

\begin{theorem}[A shared crossing beyond the first exceptional head]
\label{late}
Assume an actual extendable six-point Table 2 head at the first
three labels and rank failure, so $n=2r\ge18$ under the complete
critical-head reduction. Among all unique crossings with a shared
label to their right choose an index $i$ maximal, and let $h$ be
its last shared label. If $h\ge3$, then either the actual last three
generators form a six-point Table 2 head or $G_{r-1}$ has a singleton
orbit.
\end{theorem}
\begin{proof}
Suppose neither conclusion holds. We follow the entire acquired
maximal-crossing construction, recording all places where its former
exclusion of \emph{both} exceptional ends must be replaced.
Its interval-choice argument and the transitive/intransitive support
bounds give the actual orientation and
\[
 h<r-1,\qquad X=\operatorname{supp}(G_{>h}),\qquad
                       n-|X|\ge2h+1.                 \tag{1}
\]
If $G_{>h}$ were transitive on $X$, Lemma~\ref{location} would
force a last head: its first-head exit has $h=2$, contrary to the
hypothesis. Hence this tail is intransitive on its support.
The full ambient terminal inequality and (1) give a smallest
unique crossing $j>h$, just as in the acquired proof; otherwise
$2(r-1-h)\le |X|-2$ would contradict rank failure.

The two actual bridges yield the same regions $A,M,D$, with the
$j$-crossing separating $C=A\cup M$ from $D$. If $D$ is a
singleton, the complete interior-singleton bound makes $j=r-1$,
contrary to our exclusion. Assume $D$ nontrivial.

If this later cut is unshared, its actual sides are the supported
transitive prefix $G_{<j}$ and suffix $G_{>j}$. Here $j\ge4$,
and $j=4$ is excluded by the complete early exceptional-cut lemma:
in the critical head $G_4$ has at least three orbits. Thus $j\ge5$.
The complete later-unshared-cut theorem in
\emph{Opposite exceptional ends after a later unshared cut}
applies. Its only residual alternatives are the terminal singleton
or a degree-five dihedral suffix whose full attachment creates an
actual last head. Both contradict our exclusions. This step uses
its complete small proper-prefix exclusions and the actual natural
prefix conclusion, not raw Table 2 or a quotient IP.

There is therefore a first shared label $g$ at $j$. Maximality
of $i$ places it to the left. The entire orientation argument of the
acquired maximal-crossing proof gives
\[
 Y=\operatorname{supp}(G_{<g}),\qquad
                 n-|Y|\ge2(r-1-g)+1.                 \tag{2}
\]
Its nested-bridge count first proves $g\ge h$ and the requisite
purity of the two bridge involutions; the complete touching
exclusions then give $g>h$. None of these steps uses absence of
the first head.

If $G_{<g}$ were transitive on $Y$, apply Lemma~\ref{location}
in this proved reversed orientation. Its first-head exit requires
$g=3$, impossible because $g>h\ge3$. Its other exit is the
excluded last head. Thus both $G_{<g}$ and $G_{>h}$ are
intransitive on their actual supports.
Every colour strictly between $h$ and $g$ has two ambient-eligible
edges, by minimality of $j>h$ and $g<j$. The full orbit-incidence
inequality with (1)--(2) now gives $2r\le n-1$, the same final
contradiction as the acquired proof. These are exactly its three
exceptional exits, each replaced without assuming the already
present first head absent.
\end{proof}

\paragraph{Exact residual and source scope.}
This sharpens the either-end alternative supporting source
Proposition 5.28, physical pages 36--37, only in the displayed
critical first-head situation with $h\ge3$.
The complete early omission calculation shows that $h=2$ can occur
here only for $i=0$ in the extendable Table 2(2) orientation:
colour $1$ has two eligible edges, while the Table 2(3) initial
unique crossing has no shared labels. If any unique crossing at
index at least $5$ has a right-shared label, the maximal choice
in Theorem~\ref{late} automatically has $h\ge6$.
The first-head case with no such later right-shared crossing remains
unresolved by this note. So do elimination of the two possible
opposite ends, the stronger source rigidity/table contracts and
complete L044 recovery. No new computation, formal theorem,
original credit or question claim is made.
\end{document}
